\documentclass[8pt,a4paper]{extarticle}
% ============================================================
% Gemeinsamer Preamble fuer alle 5 Open-Book-Spickzettel.
% Wird aus jedem klausur.tex via % ============================================================
% Gemeinsamer Preamble fuer alle 5 Open-Book-Spickzettel.
% Wird aus jedem klausur.tex via % ============================================================
% Gemeinsamer Preamble fuer alle 5 Open-Book-Spickzettel.
% Wird aus jedem klausur.tex via \input{../preamble.tex} geladen.
% Layout: A4 hochkant, 3 Spalten, 8pt, sehr kompakt.
% ============================================================

\usepackage[a4paper,margin=8mm,top=7mm,bottom=7mm,includefoot]{geometry}
\usepackage[utf8]{inputenc}
\usepackage[T1]{fontenc}
\usepackage[ngerman]{babel}
\usepackage{lmodern}
\usepackage{microtype}

\usepackage{amsmath,amssymb,amsfonts,mathtools}
\usepackage{multicol}
\usepackage{enumitem}
\usepackage{array,tabularx,booktabs,makecell}
\usepackage[dvipsnames,table]{xcolor}
\usepackage[explicit]{titlesec}
\usepackage{etoolbox}
\usepackage{fancyhdr}
\usepackage{lastpage}

% ---------- Farben ----------
\definecolor{secblue}{HTML}{1f4e79}
\definecolor{subblue}{HTML}{2e75b6}
\definecolor{boxbg}{HTML}{f2f2f2}
\definecolor{accent}{HTML}{c00000}
\definecolor{rulegray}{HTML}{888888}

% ---------- Spalten ----------
\setlength{\columnsep}{3mm}
\setlength{\columnseprule}{0.3pt}
\def\columnseprulecolor{\color{rulegray}}

% ---------- Listen kompakt ----------
\setlist{nosep,leftmargin=1.1em,labelsep=0.35em,topsep=0.2ex,partopsep=0pt}
\setlist[itemize]{label=\textbullet}

% ---------- Abstaende ----------
\setlength{\parindent}{0pt}
\setlength{\parskip}{0.15ex}
\setlength{\abovedisplayskip}{0.3ex}
\setlength{\belowdisplayskip}{0.3ex}
\setlength{\abovedisplayshortskip}{0.2ex}
\setlength{\belowdisplayshortskip}{0.2ex}
\setlength{\jot}{1pt}

% ---------- Tabellen kompakt ----------
\renewcommand{\arraystretch}{1.05}
\setlength{\tabcolsep}{3pt}
\setlength{\arrayrulewidth}{0.3pt}
\arrayrulecolor{rulegray}

% ---------- Section-Stil (kein Vollblock; dick + farbige Linie darunter) ----------
\titleformat{\section}
  {\large\bfseries\color{secblue}}
  {}{0pt}{#1}
  [\vspace{-0.4ex}{\color{secblue}\hrule height 0.8pt}\vspace{0.1ex}]
\titlespacing*{\section}{0pt}{1.4ex plus 0.4ex}{0.5ex}

\titleformat{\subsection}
  {\small\bfseries\color{subblue}}
  {}{0pt}{#1}
\titlespacing*{\subsection}{0pt}{0.9ex plus 0.2ex}{0.2ex}

\titleformat{\subsubsection}
  {\small\bfseries}
  {}{0pt}{#1}
\titlespacing*{\subsubsection}{0pt}{0.5ex}{0.1ex}

% ---------- Mini-Helfer ----------

% \fact{...} Hervorgehobene Regel (gerahmt, hellgrau)
\newcommand{\fact}[1]{%
  \begingroup\setlength{\fboxsep}{2pt}%
  \colorbox{boxbg}{\parbox{\dimexpr\linewidth-2\fboxsep}{\small #1\strut}}%
  \endgroup\par\vspace{0.2ex}}

% \warn{...} Stolperfalle (roter Strich links)
\newcommand{\warn}[1]{%
  \begingroup\setlength{\fboxsep}{2pt}%
  \noindent{\color{accent}\rule{2pt}{\baselineskip}}\hspace{2pt}%
  \begin{minipage}[t]{\dimexpr\linewidth-6pt}\small #1\end{minipage}%
  \endgroup\par\vspace{0.2ex}}

% Math-Shortcuts
\newcommand{\R}{\mathbb{R}}
\newcommand{\N}{\mathbb{N}}
\newcommand{\Z}{\mathbb{Z}}
\newcommand{\Q}{\mathbb{Q}}
\newcommand{\C}{\mathbb{C}}
\DeclareMathOperator{\rg}{rg}
\DeclareMathOperator{\spur}{tr}
\DeclareMathOperator{\im}{im}
\DeclareMathOperator{\diag}{diag}
\DeclareMathOperator{\sign}{sgn}
\DeclareMathOperator*{\argmin}{arg\,min}
\DeclareMathOperator*{\argmax}{arg\,max}
\newcommand{\trans}{^{\!\top}}
\newcommand{\norm}[1]{\left\lVert #1 \right\rVert}
\newcommand{\abs}[1]{\left\lvert #1 \right\rvert}
\newcommand{\inner}[2]{\langle #1,\,#2\rangle}
\newcommand{\dx}{\,\mathrm{d}x}
\newcommand{\dy}{\,\mathrm{d}y}
\newcommand{\dt}{\,\mathrm{d}t}

% Zeigt eine Display-Formel zentriert; skaliert nur runter wenn zu breit.
\newlength{\fitw}
\newcommand{\fitline}[1]{%
  \par\nointerlineskip
  \settowidth{\fitw}{$\displaystyle #1$}%
  \ifdim\fitw>\linewidth
    \noindent\resizebox{\linewidth}{!}{$\displaystyle #1$}\par
  \else
    \[\displaystyle #1\]%
  \fi}

% Tabular fuer fixe Spaltenbreiten (Spickzettel-typisch)
\newenvironment{tightt}[1]
  {\par\noindent\begin{tabular}{#1}}
  {\end{tabular}\par}
% Fuer X-Spalten direkt \begin{tabularx}{\linewidth}{...} verwenden
% (kann nicht in \newenvironment gewrappt werden -- TX@get@body Konflikt)

% Globale Schutzschalter gegen Overflow
\emergencystretch=2em
\tolerance=2000
\hbadness=10000
\hfuzz=2pt
\vfuzz=2pt

% ---------- Kopf-/Fusszeile ----------
\pagestyle{fancy}
\fancyhf{}
\renewcommand{\headrulewidth}{0pt}
\renewcommand{\footrulewidth}{0pt}
\fancyfoot[L]{\scriptsize\color{rulegray}\textit{\sheettitle}}
\fancyfoot[R]{\scriptsize\color{rulegray}Seite \thepage\,/\,\pageref{LastPage}}
\setlength{\footskip}{8pt}

% Titel-Header oben auf Seite 1
\newcommand{\sheetheader}[2]{%
  {\Large\bfseries\color{secblue} #1}\hfill
  {\small\color{rulegray} #2}\par
  \vspace{0.4ex}{\color{secblue}\hrule height 0.7pt}\vspace{1ex}}

% Default-Titel (wird ueberschrieben)
\newcommand{\sheettitle}{Open-Book Spickzettel}
 geladen.
% Layout: A4 hochkant, 3 Spalten, 8pt, sehr kompakt.
% ============================================================

\usepackage[a4paper,margin=8mm,top=7mm,bottom=7mm,includefoot]{geometry}
\usepackage[utf8]{inputenc}
\usepackage[T1]{fontenc}
\usepackage[ngerman]{babel}
\usepackage{lmodern}
\usepackage{microtype}

\usepackage{amsmath,amssymb,amsfonts,mathtools}
\usepackage{multicol}
\usepackage{enumitem}
\usepackage{array,tabularx,booktabs,makecell}
\usepackage[dvipsnames,table]{xcolor}
\usepackage[explicit]{titlesec}
\usepackage{etoolbox}
\usepackage{fancyhdr}
\usepackage{lastpage}

% ---------- Farben ----------
\definecolor{secblue}{HTML}{1f4e79}
\definecolor{subblue}{HTML}{2e75b6}
\definecolor{boxbg}{HTML}{f2f2f2}
\definecolor{accent}{HTML}{c00000}
\definecolor{rulegray}{HTML}{888888}

% ---------- Spalten ----------
\setlength{\columnsep}{3mm}
\setlength{\columnseprule}{0.3pt}
\def\columnseprulecolor{\color{rulegray}}

% ---------- Listen kompakt ----------
\setlist{nosep,leftmargin=1.1em,labelsep=0.35em,topsep=0.2ex,partopsep=0pt}
\setlist[itemize]{label=\textbullet}

% ---------- Abstaende ----------
\setlength{\parindent}{0pt}
\setlength{\parskip}{0.15ex}
\setlength{\abovedisplayskip}{0.3ex}
\setlength{\belowdisplayskip}{0.3ex}
\setlength{\abovedisplayshortskip}{0.2ex}
\setlength{\belowdisplayshortskip}{0.2ex}
\setlength{\jot}{1pt}

% ---------- Tabellen kompakt ----------
\renewcommand{\arraystretch}{1.05}
\setlength{\tabcolsep}{3pt}
\setlength{\arrayrulewidth}{0.3pt}
\arrayrulecolor{rulegray}

% ---------- Section-Stil (kein Vollblock; dick + farbige Linie darunter) ----------
\titleformat{\section}
  {\large\bfseries\color{secblue}}
  {}{0pt}{#1}
  [\vspace{-0.4ex}{\color{secblue}\hrule height 0.8pt}\vspace{0.1ex}]
\titlespacing*{\section}{0pt}{1.4ex plus 0.4ex}{0.5ex}

\titleformat{\subsection}
  {\small\bfseries\color{subblue}}
  {}{0pt}{#1}
\titlespacing*{\subsection}{0pt}{0.9ex plus 0.2ex}{0.2ex}

\titleformat{\subsubsection}
  {\small\bfseries}
  {}{0pt}{#1}
\titlespacing*{\subsubsection}{0pt}{0.5ex}{0.1ex}

% ---------- Mini-Helfer ----------

% \fact{...} Hervorgehobene Regel (gerahmt, hellgrau)
\newcommand{\fact}[1]{%
  \begingroup\setlength{\fboxsep}{2pt}%
  \colorbox{boxbg}{\parbox{\dimexpr\linewidth-2\fboxsep}{\small #1\strut}}%
  \endgroup\par\vspace{0.2ex}}

% \warn{...} Stolperfalle (roter Strich links)
\newcommand{\warn}[1]{%
  \begingroup\setlength{\fboxsep}{2pt}%
  \noindent{\color{accent}\rule{2pt}{\baselineskip}}\hspace{2pt}%
  \begin{minipage}[t]{\dimexpr\linewidth-6pt}\small #1\end{minipage}%
  \endgroup\par\vspace{0.2ex}}

% Math-Shortcuts
\newcommand{\R}{\mathbb{R}}
\newcommand{\N}{\mathbb{N}}
\newcommand{\Z}{\mathbb{Z}}
\newcommand{\Q}{\mathbb{Q}}
\newcommand{\C}{\mathbb{C}}
\DeclareMathOperator{\rg}{rg}
\DeclareMathOperator{\spur}{tr}
\DeclareMathOperator{\im}{im}
\DeclareMathOperator{\diag}{diag}
\DeclareMathOperator{\sign}{sgn}
\DeclareMathOperator*{\argmin}{arg\,min}
\DeclareMathOperator*{\argmax}{arg\,max}
\newcommand{\trans}{^{\!\top}}
\newcommand{\norm}[1]{\left\lVert #1 \right\rVert}
\newcommand{\abs}[1]{\left\lvert #1 \right\rvert}
\newcommand{\inner}[2]{\langle #1,\,#2\rangle}
\newcommand{\dx}{\,\mathrm{d}x}
\newcommand{\dy}{\,\mathrm{d}y}
\newcommand{\dt}{\,\mathrm{d}t}

% Zeigt eine Display-Formel zentriert; skaliert nur runter wenn zu breit.
\newlength{\fitw}
\newcommand{\fitline}[1]{%
  \par\nointerlineskip
  \settowidth{\fitw}{$\displaystyle #1$}%
  \ifdim\fitw>\linewidth
    \noindent\resizebox{\linewidth}{!}{$\displaystyle #1$}\par
  \else
    \[\displaystyle #1\]%
  \fi}

% Tabular fuer fixe Spaltenbreiten (Spickzettel-typisch)
\newenvironment{tightt}[1]
  {\par\noindent\begin{tabular}{#1}}
  {\end{tabular}\par}
% Fuer X-Spalten direkt \begin{tabularx}{\linewidth}{...} verwenden
% (kann nicht in \newenvironment gewrappt werden -- TX@get@body Konflikt)

% Globale Schutzschalter gegen Overflow
\emergencystretch=2em
\tolerance=2000
\hbadness=10000
\hfuzz=2pt
\vfuzz=2pt

% ---------- Kopf-/Fusszeile ----------
\pagestyle{fancy}
\fancyhf{}
\renewcommand{\headrulewidth}{0pt}
\renewcommand{\footrulewidth}{0pt}
\fancyfoot[L]{\scriptsize\color{rulegray}\textit{\sheettitle}}
\fancyfoot[R]{\scriptsize\color{rulegray}Seite \thepage\,/\,\pageref{LastPage}}
\setlength{\footskip}{8pt}

% Titel-Header oben auf Seite 1
\newcommand{\sheetheader}[2]{%
  {\Large\bfseries\color{secblue} #1}\hfill
  {\small\color{rulegray} #2}\par
  \vspace{0.4ex}{\color{secblue}\hrule height 0.7pt}\vspace{1ex}}

% Default-Titel (wird ueberschrieben)
\newcommand{\sheettitle}{Open-Book Spickzettel}
 geladen.
% Layout: A4 hochkant, 3 Spalten, 8pt, sehr kompakt.
% ============================================================

\usepackage[a4paper,margin=8mm,top=7mm,bottom=7mm,includefoot]{geometry}
\usepackage[utf8]{inputenc}
\usepackage[T1]{fontenc}
\usepackage[ngerman]{babel}
\usepackage{lmodern}
\usepackage{microtype}

\usepackage{amsmath,amssymb,amsfonts,mathtools}
\usepackage{multicol}
\usepackage{enumitem}
\usepackage{array,tabularx,booktabs,makecell}
\usepackage[dvipsnames,table]{xcolor}
\usepackage[explicit]{titlesec}
\usepackage{etoolbox}
\usepackage{fancyhdr}
\usepackage{lastpage}

% ---------- Farben ----------
\definecolor{secblue}{HTML}{1f4e79}
\definecolor{subblue}{HTML}{2e75b6}
\definecolor{boxbg}{HTML}{f2f2f2}
\definecolor{accent}{HTML}{c00000}
\definecolor{rulegray}{HTML}{888888}

% ---------- Spalten ----------
\setlength{\columnsep}{3mm}
\setlength{\columnseprule}{0.3pt}
\def\columnseprulecolor{\color{rulegray}}

% ---------- Listen kompakt ----------
\setlist{nosep,leftmargin=1.1em,labelsep=0.35em,topsep=0.2ex,partopsep=0pt}
\setlist[itemize]{label=\textbullet}

% ---------- Abstaende ----------
\setlength{\parindent}{0pt}
\setlength{\parskip}{0.15ex}
\setlength{\abovedisplayskip}{0.3ex}
\setlength{\belowdisplayskip}{0.3ex}
\setlength{\abovedisplayshortskip}{0.2ex}
\setlength{\belowdisplayshortskip}{0.2ex}
\setlength{\jot}{1pt}

% ---------- Tabellen kompakt ----------
\renewcommand{\arraystretch}{1.05}
\setlength{\tabcolsep}{3pt}
\setlength{\arrayrulewidth}{0.3pt}
\arrayrulecolor{rulegray}

% ---------- Section-Stil (kein Vollblock; dick + farbige Linie darunter) ----------
\titleformat{\section}
  {\large\bfseries\color{secblue}}
  {}{0pt}{#1}
  [\vspace{-0.4ex}{\color{secblue}\hrule height 0.8pt}\vspace{0.1ex}]
\titlespacing*{\section}{0pt}{1.4ex plus 0.4ex}{0.5ex}

\titleformat{\subsection}
  {\small\bfseries\color{subblue}}
  {}{0pt}{#1}
\titlespacing*{\subsection}{0pt}{0.9ex plus 0.2ex}{0.2ex}

\titleformat{\subsubsection}
  {\small\bfseries}
  {}{0pt}{#1}
\titlespacing*{\subsubsection}{0pt}{0.5ex}{0.1ex}

% ---------- Mini-Helfer ----------

% \fact{...} Hervorgehobene Regel (gerahmt, hellgrau)
\newcommand{\fact}[1]{%
  \begingroup\setlength{\fboxsep}{2pt}%
  \colorbox{boxbg}{\parbox{\dimexpr\linewidth-2\fboxsep}{\small #1\strut}}%
  \endgroup\par\vspace{0.2ex}}

% \warn{...} Stolperfalle (roter Strich links)
\newcommand{\warn}[1]{%
  \begingroup\setlength{\fboxsep}{2pt}%
  \noindent{\color{accent}\rule{2pt}{\baselineskip}}\hspace{2pt}%
  \begin{minipage}[t]{\dimexpr\linewidth-6pt}\small #1\end{minipage}%
  \endgroup\par\vspace{0.2ex}}

% Math-Shortcuts
\newcommand{\R}{\mathbb{R}}
\newcommand{\N}{\mathbb{N}}
\newcommand{\Z}{\mathbb{Z}}
\newcommand{\Q}{\mathbb{Q}}
\newcommand{\C}{\mathbb{C}}
\DeclareMathOperator{\rg}{rg}
\DeclareMathOperator{\spur}{tr}
\DeclareMathOperator{\im}{im}
\DeclareMathOperator{\diag}{diag}
\DeclareMathOperator{\sign}{sgn}
\DeclareMathOperator*{\argmin}{arg\,min}
\DeclareMathOperator*{\argmax}{arg\,max}
\newcommand{\trans}{^{\!\top}}
\newcommand{\norm}[1]{\left\lVert #1 \right\rVert}
\newcommand{\abs}[1]{\left\lvert #1 \right\rvert}
\newcommand{\inner}[2]{\langle #1,\,#2\rangle}
\newcommand{\dx}{\,\mathrm{d}x}
\newcommand{\dy}{\,\mathrm{d}y}
\newcommand{\dt}{\,\mathrm{d}t}

% Zeigt eine Display-Formel zentriert; skaliert nur runter wenn zu breit.
\newlength{\fitw}
\newcommand{\fitline}[1]{%
  \par\nointerlineskip
  \settowidth{\fitw}{$\displaystyle #1$}%
  \ifdim\fitw>\linewidth
    \noindent\resizebox{\linewidth}{!}{$\displaystyle #1$}\par
  \else
    \[\displaystyle #1\]%
  \fi}

% Tabular fuer fixe Spaltenbreiten (Spickzettel-typisch)
\newenvironment{tightt}[1]
  {\par\noindent\begin{tabular}{#1}}
  {\end{tabular}\par}
% Fuer X-Spalten direkt \begin{tabularx}{\linewidth}{...} verwenden
% (kann nicht in \newenvironment gewrappt werden -- TX@get@body Konflikt)

% Globale Schutzschalter gegen Overflow
\emergencystretch=2em
\tolerance=2000
\hbadness=10000
\hfuzz=2pt
\vfuzz=2pt

% ---------- Kopf-/Fusszeile ----------
\pagestyle{fancy}
\fancyhf{}
\renewcommand{\headrulewidth}{0pt}
\renewcommand{\footrulewidth}{0pt}
\fancyfoot[L]{\scriptsize\color{rulegray}\textit{\sheettitle}}
\fancyfoot[R]{\scriptsize\color{rulegray}Seite \thepage\,/\,\pageref{LastPage}}
\setlength{\footskip}{8pt}

% Titel-Header oben auf Seite 1
\newcommand{\sheetheader}[2]{%
  {\Large\bfseries\color{secblue} #1}\hfill
  {\small\color{rulegray} #2}\par
  \vspace{0.4ex}{\color{secblue}\hrule height 0.7pt}\vspace{1ex}}

% Default-Titel (wird ueberschrieben)
\newcommand{\sheettitle}{Open-Book Spickzettel}

\renewcommand{\sheettitle}{Analysis 2}

\begin{document}
\sheetheader{Analysis 2}{Open-Book Formelsammlung -- DHBW WDS125}

\begin{multicols*}{3}

% =================================================
\section{Folgen \& Reihen}
% =================================================
\subsection{Folgen}
Konv.: $\forall\epsilon\!>\!0\,\exists N\!:\,\forall n\!>\!N\!:\,|a_n-a|<\epsilon$.\\
\textbf{Standardgrenzwerte}:
\begin{itemize}
\item $\lim 1/n^p=0$ ($p>0$)
\item $\lim\sqrt[n]{n}=1$,\ $\lim\sqrt[n]{a}=1$ ($a>0$)
\item $\lim(1+1/n)^n=e$
\item $\lim q^n=0$ ($|q|<1$)
\end{itemize}
\textbf{Sandwich}: $a_n\!\le\!b_n\!\le\!c_n$, $a_n,c_n\!\to\!L$ $\Rightarrow$ $b_n\!\to\!L$.\\
\textbf{Monoton + beschr\"ankt} $\Rightarrow$ konvergent.

\subsection{Reihen}
$\sum a_n$ konv.\ $\Leftrightarrow$ Partialsummen $s_n$ konv.
\fact{\textbf{Notwendig}: $\sum a_n$ konv.\ $\Rightarrow a_n\!\to\!0$ (\emph{nicht} hinreichend, vgl.\ harmonische).}

\subsection{Konvergenzkriterien}
\begin{tabularx}{\linewidth}{@{}l@{\ }X@{}}
Geom.\ $\sum q^n$ & konv.\ $\Leftrightarrow$ $|q|<1$, Wert $1/(1{-}q)$ \\
Majorante & $0\!\le\!|a_n|\!\le\!b_n,\,\sum b_n$ konv.\ $\Rightarrow$ abs.\ konv. \\
Minorante & $a_n\!\ge\!b_n\!\ge\!0,\,\sum b_n$ div.\ $\Rightarrow$ div. \\
Quotient & $\lim|a_{n+1}/a_n|{=}q$: $q\!<\!1$ konv., $q\!>\!1$ div. \\
Wurzel & $\lim\sqrt[n]{|a_n|}{=}q$ analog \\
Leibniz & altern., $|a_n|\!\downarrow,\!\to\!0$ $\Rightarrow$ konv. \\
Integralvgl. & $f$ pos.\ \& fallend: $\sum f(n)\sim\!\int f$ \\
Verdichtung & $a_n\!\downarrow\!>0$: $\sum a_n\sim\sum 2^n a_{2^n}$ \\
\end{tabularx}
\textbf{Absolut konv.} ($\sum|a_n|$ konv.) erlaubt Umordnung. \textbf{Bedingt}: konv.\ aber nicht absolut (Riemannscher Umordnungssatz).

\subsection{Standardreihen}
\begin{tabularx}{\linewidth}{@{}l@{\ }X@{}}
$\sum 1/n^s$ & konv.\ $\Leftrightarrow s>1$;\ $\zeta(2){=}\pi^2/6$ \\
$\sum(-1)^{n+1}/n$ & $=\ln 2$ (Leibniz) \\
$\sum 1/n!$ & $=e-1$ (ab $n{=}1$) \\
$\sum q^n$ & $=1/(1-q)$, $|q|<1$ \\
\end{tabularx}

\subsection{Potenzreihen}
$\sum_{n=0}^\infty a_n(x-x_0)^n$,\ Radius $r=1/\limsup\sqrt[n]{|a_n|}=\lim|a_n/a_{n+1}|$.\\
Innerhalb gliedweise diff./integr.

\subsection{Taylor}
\fitline{f(x)=\sum_{n=0}^\infty \tfrac{f^{(n)}(x_0)}{n!}(x-x_0)^n}
\textbf{Maclaurin}:
\begin{itemize}
\item $e^x=\sum x^n/n!$,\ $r=\infty$
\item $\sin x=\sum (-1)^n x^{2n+1}/(2n{+}1)!$,\ $r=\infty$
\item $\cos x=\sum (-1)^n x^{2n}/(2n)!$,\ $r=\infty$
\item $\ln(1{+}x)=\sum_{n\ge 1}(-1)^{n+1}x^n/n$,\ $r=1$
\item $(1{+}x)^\alpha=\sum\binom{\alpha}{n}x^n$,\ $r=1$
\item $1/(1{-}x)=\sum x^n$,\ $r=1$
\end{itemize}
\textbf{Restglied (Lagrange)}: $R_n(x)=\tfrac{f^{(n+1)}(\xi)}{(n+1)!}(x{-}x_0)^{n+1}$.

% =================================================
\section{Mehrdim.\ Differentialrechnung}
% =================================================
\subsection{Topologie \& Stetigkeit}
$\|x\|=\sqrt{\sum x_i^2}$;\ $B_r(a)=\{x:\|x-a\|<r\}$.\\
$M$ offen $\Leftrightarrow$ jedes $a\!\in\!M$ hat $B_r(a)\!\subseteq\!M$.\\
$f$ stetig: $\lim_{x\to a}f(x)=f(a)$.

\subsection{Gradient}
\fitline{\nabla f(a)=(\partial_{x_1}f,\dots,\partial_{x_n}f)\trans}
Steilster Anstieg;\ $\|\nabla f\|$ = max.\ Steigung;\ $\nabla f\perp$ Niveaulinie.

\subsection{Totale Differenzierbarkeit}
$f$ total diff.\ in $a$: $\exists L$ lin.\ mit $f(a{+}h)=f(a)+Lh+o(\|h\|)$.
\warn{Partielle Ableitungen existent $\ne$ differenzierbar. Hinreichend: alle partiellen Ableitungen in Umgebung stetig.}

\subsection{Richtungsableitung}
$D_v f(a)=\lim_{t\to 0}(f(a{+}tv){-}f(a))/t$.\\
Bei total diff.: $D_v f(a)=\nabla f(a)\trans v$.\\
Maximal bei $v=\nabla f/\|\nabla f\|$.

\subsection{Kettenregel}
\fitline{J_{g\circ f}(x)=J_g(f(x))\cdot J_f(x)}
Kurve: $h(t)=f(\gamma(t))$, $h'(t)=\nabla f(\gamma(t))\trans \gamma'(t)$.

\subsection{Hesse \& Taylor}
$H_f(x)=(\partial^2_{x_i x_j}f)_{ij}$.
\fact{\textbf{Schwarz}: stetige 2.\ Ableitungen $\Rightarrow$ $H$ symmetrisch.}
\fitline{f(a{+}h)=f(a)+\nabla f(a)\trans h+\tfrac12 h\trans H_f(a) h+R_2}

% =================================================
\section{Extrema \& Lagrange}
% =================================================
\subsection{Freie Extrema}
\begin{enumerate}
\item $\nabla f(x){=}0$ $\Rightarrow$ station\"are Punkte
\item $H_f(x^\star)$ bestimmen
\item Klassifizieren via EW:
\end{enumerate}
\begin{tabularx}{\linewidth}{@{}l@{\ }X@{}}
PD (alle EW $>0$) & striktes lok.\ Min \\
ND (alle EW $<0$) & striktes lok.\ Max \\
indefinit & Sattel \\
semidef.\ (EW $=0$) & unentschieden \\
\end{tabularx}
\textbf{2$\times$2-Schnelltest} $H=\bigl(\begin{smallmatrix}A&B\\B&C\end{smallmatrix}\bigr)$, $\det H{=}AC{-}B^2$:
\begin{itemize}
\item $\det H>0,\ A>0$ $\Rightarrow$ Min
\item $\det H>0,\ A<0$ $\Rightarrow$ Max
\item $\det H<0$ $\Rightarrow$ Sattel
\item $\det H=0$ $\Rightarrow$ unklar
\end{itemize}

\subsection{Lagrange-Multiplikatoren (Gleichungs-NB)}
$\min f(x)$ u.d.N.\ $g_i(x){=}0$:
\fitline{L(x,\lambda)=f(x)-\sum_i \lambda_i g_i(x)}
Notw.: $\nabla f(x^\star)=\sum_i \lambda_i \nabla g_i(x^\star)$.

% =================================================
\section{Vektoranalysis}
% =================================================
F\"ur $F=(P,Q,R)$:
\fitline{\nabla\cdot F=\partial_x P+\partial_y Q+\partial_z R\ \text{(Quellendichte)}}
\fitline{\nabla\times F=\bigl(\partial_y R{-}\partial_z Q,\,\partial_z P{-}\partial_x R,\,\partial_x Q{-}\partial_y P\bigr)}
\textbf{Identit\"aten}:
\begin{itemize}
\item $\nabla\times(\nabla f)=0$
\item $\nabla\cdot(\nabla\times F)=0$
\item $\nabla\times(\nabla\times F)=\nabla(\nabla\cdot F)-\Delta F$
\end{itemize}

\subsection{Kurvenintegrale}
\fitline{\textstyle\int_\gamma F\cdot dr=\!\int_a^b\! F(\gamma(t))\trans\gamma'(t)\,dt}
\textbf{Konservativ}: $F=\nabla\phi\Rightarrow$ wegunabh.; $\int_\gamma F\cdot dr=\phi(\gamma(b))-\phi(\gamma(a))$.\\
Einfach zsh.\ Gebiet: konservativ $\Leftrightarrow \nabla\times F=0$.

\subsection{Integrals\"atze}
\textbf{Gauss}: $\iiint_V\!\nabla\!\cdot\!F\,dV=\oiint_{\partial V}\!F\!\cdot\!dS$\\
\textbf{Stokes}: $\iint_S\!(\nabla\!\times\!F)\!\cdot\!dS=\oint_{\partial S}\!F\!\cdot\!dr$\\
\textbf{Green} (2D): $\iint_D(Q_x{-}P_y)\,dA=\oint_{\partial D}\!P\,dx+Q\,dy$

% =================================================
\section{Mehrfachintegrale}
% =================================================
\subsection{Normalbereich}
\fitline{\iint_D f\,dA=\!\int_a^b\!\int_{g_1(x)}^{g_2(x)}\!f\,dy\,dx}
\textbf{Fubini}: Reihenfolge tauschbar bei genuegend Regularit\"at.

\subsection{Transformationssatz}
\fitline{\iint_D f(x,y)\,dA=\iint_{D'}\!f(x(u,v),y(u,v))\,|\det J|\,du\,dv}
\begin{tabularx}{\linewidth}{@{}l@{\ }X@{}}
Polar (2D) & $x{=}r\cos\phi,\ y{=}r\sin\phi$;\ $|J|{=}r$ \\
Zylinder & wie Polar, $z{=}z$;\ $|J|{=}r$ \\
Kugel & $x{=}r\sin\theta\cos\phi,\ y{=}r\sin\theta\sin\phi,\ z{=}r\cos\theta$;\ $|J|{=}r^2\sin\theta$ \\
\end{tabularx}
\fact{$\int_{-\infty}^\infty e^{-x^2}\,dx=\sqrt\pi$;\ Kugelvol.\ $V{=}\tfrac43\pi R^3$.}

% =================================================
\section{ODE -- Klassifikation \& Loesungen}
% =================================================
Ordnung = h\"ochste Ableitung;\ linear: $\sum a_k(x)y^{(k)}=b(x)$;\ homogen: $b{=}0$;\ konst.\ Koeff.: $a_k$ konst.\\
AWP = ODE + Werte $y(x_0),y'(x_0)$;\ RWP = Werte an beiden Endpunkten.

\subsection{ODE 1.\ Ordnung}
\textbf{Trennbar} $y'=f(x)g(y)$:
\fitline{\textstyle\int\!\tfrac{dy}{g(y)}=\!\int\!f(x)\,dx+C}
\textbf{Linear} $y'+p(x)y=q(x)$, $\mu(x)=e^{\int p\,dx}$:
\fitline{y=\tfrac{1}{\mu}\!\left(\!\textstyle\int\!\mu(x)q(x)\,dx+C\right)}
\textbf{Bernoulli} $y'+py=qy^n,\ n\!\ne\!1$: Subst.\ $u=y^{1-n}\Rightarrow$ linear in $u$.\\
\textbf{Exakt} $M\,dx+N\,dy=0,\ M_y=N_x$: $\Phi$ mit $\Phi_x{=}M,\Phi_y{=}N$;\ L\"osung $\Phi(x,y)=C$.\\
\textbf{Homogen Grad 0}: $y'=f(y/x)\Rightarrow u{=}y/x$.\\
\textbf{Riccati} $y'=a+by+cy^2$ mit bekanntem $y_1$: $y=y_1+1/u\Rightarrow$ linear.

\subsection{ODE 2.\ Ord., konst.\ Koeff.}
$y''+ay'+by=f(x)$. Char.\ Gl.\ $\lambda^2+a\lambda+b=0$:
\begin{tabularx}{\linewidth}{@{}l@{\ }X@{}}
$a^2{-}4b>0$ & $\lambda_1\!\ne\!\lambda_2$: $C_1 e^{\lambda_1 x}{+}C_2 e^{\lambda_2 x}$ \\
$=0$ & $\lambda$ doppelt: $(C_1{+}C_2 x)e^{\lambda x}$ \\
$<0$ & $\alpha{\pm}i\beta$: $e^{\alpha x}(C_1\!\cos\!\beta x{+}C_2\!\sin\!\beta x)$ \\
\end{tabularx}
\textbf{Inhomogen -- Ansatz nach $f$}:
\begin{tabularx}{\linewidth}{@{}l@{\ }X@{}}
Polynom Grad $n$ & Polynom Grad $n$ \\
$e^{kx}$ & $A e^{kx}$ \\
$\cos\omega x,\sin\omega x$ & $A\cos\omega x+B\sin\omega x$ \\
Produkt & Produkt der Ans\"atze \\
\end{tabularx}
\fact{\textbf{Resonanz}: $f$ ist L\"osung der homogenen $\Rightarrow$ Ansatz $\cdot x$ (doppelt: $\cdot x^2$).}
\textbf{Variation der Konst.}:
\fitline{y_p=-y_1\!\int\!\tfrac{y_2 f}{W}\,dx+y_2\!\int\!\tfrac{y_1 f}{W}\,dx,\ W{=}y_1 y_2'{-}y_2 y_1'}

\subsection{Systeme linearer ODE}
$\mathbf y'=A\mathbf y,\ \mathbf y(0)=\mathbf y_0$:
\fitline{\mathbf y(t)=e^{At}\mathbf y_0,\ e^{At}=\sum_{k\ge 0}\tfrac{(At)^k}{k!}}
Diag.\ $A=PDP^{-1}$: $e^{At}=Pe^{Dt}P^{-1}$, $e^{Dt}=\diag(e^{\lambda_i t})$.

% =================================================
\section{Klausur-Stolperfallen}
% =================================================
\begin{itemize}
\item Notw.\ Krit.\ ($a_n\!\to\!0$) ist \emph{nicht} hinreichend; harmonisch divergiert.
\item Bei altern.\ Reihen \textbf{Leibniz}, nicht Quotient.
\item Absolut vs.\ bedingt unterscheiden.
\item Partielle Abl.\ vorhanden $\ne$ stetig/diff.bar.
\item \textbf{Resonanz} im Ansatz $\cdot x$ pr\"ufen, sonst falsch.
\item Quadrik/Hesse aus EW-\textbf{Signatur}, nicht aus $\det$ allein.
\item Schwarz nur bei stetigen 2.\ Ableitungen anwenden.
\item Transformationssatz: $|\det J|$ nicht vergessen (z.B.\ $r$ in Polar).
\item Lagrange-Multiplikator: Vorzeichen in $L=f-\lambda g$ (Konvention beachten).
\end{itemize}

\end{multicols*}
\end{document}
